\section*{Appendix}
\addcontentsline{toc}{section}{Appendix}
\label{app:appendix}
%%%%%%%%%%%%%%%%%%%%%

\subsection*{Code Repository}

The project code can be found at the following link: \href{https://github.com/vbhaga89/MIT805_Project}{GitHub Repository}

\subsection*{Spark Execution Evidence}

The following spark execution plan is for the payment format aggregation results. The \texttt{mode="formatted"} parameter was used for the \texttt{explain()} method because this provides not only the physical plan outline but also detailed information for each specific node. \autoref{code_snippet_6} illustrates the formatted execution plan for the payment format aggregation:

\begin{lstlisting}[
    caption={PySpark Execution Plan},
    label={code_snippet_6}
]
== Physical Plan ==
AdaptiveSparkPlan (9)
+- Sort (8)
   +- Exchange (7)
      +- Project (6)
         +- Project (5)
            +- HashAggregate (4)
               +- Exchange (3)
                  +- HashAggregate (2)
                     +- Scan parquet  (1)


(1) Scan parquet 
Output [2]: [Payment Format#9, Is Laundering#10L]
Batched: true
Location: InMemoryFileIndex [file:/content/drive/MyDrive/MIT805_Project/filtered_cleaned_transactions.parquet]
ReadSchema: struct<Payment Format:string,Is Laundering:bigint>

(2) HashAggregate
Input [2]: [Payment Format#9, Is Laundering#10L]
Keys [1]: [Payment Format#9]
Functions [2]: [partial_count(1), partial_sum(Is Laundering#10L)]
Aggregate Attributes [2]: [count#123L, sum#124L]
Results [3]: [Payment Format#9, count#125L, sum#126L]

(3) Exchange
Input [3]: [Payment Format#9, count#125L, sum#126L]
Arguments: hashpartitioning(Payment Format#9, 32), ENSURE_REQUIREMENTS, [plan_id=1067]

(4) HashAggregate
Input [3]: [Payment Format#9, count#125L, sum#126L]
Keys [1]: [Payment Format#9]
Functions [2]: [count(1), sum(Is Laundering#10L)]
Aggregate Attributes [2]: [count(1)#112L, sum(Is Laundering#10L)#113L]
Results [3]: [Payment Format#9, count(1)#112L AS Transactions#96L, sum(Is Laundering#10L)#113L AS Laundering#97L]

(5) Project
Output [5]: [Payment Format#9, Transactions#96L, Laundering#97L, ((cast(Laundering#97L as double) / cast(Transactions#96L as double)) * 100.0) AS Laundering Rate (%)#114, ((cast(Laundering#97L as double) / 96938.0) * 100.0) AS Share of Laundering (%)#115]
Input [3]: [Payment Format#9, Transactions#96L, Laundering#97L]
\end{lstlisting}

\newpage

\begin{lstlisting}[
    title={PySpark Execution Plan (continued)}
]
(6) Project
Output [6]: [Payment Format#9, Transactions#96L, Laundering#97L, Laundering Rate (%)#114, Share of Laundering (%)#115, (Laundering Rate (%)#114 / 0.055059604671720386) AS Relative Risk#116]
Input [5]: [Payment Format#9, Transactions#96L, Laundering#97L, Laundering Rate (%)#114, Share of Laundering (%)#115]

(7) Exchange
Input [6]: [Payment Format#9, Transactions#96L, Laundering#97L, Laundering Rate (%)#114, Share of Laundering (%)#115, Relative Risk#116]
Arguments: rangepartitioning(Laundering Rate (%)#114 DESC NULLS LAST, 32), ENSURE_REQUIREMENTS, [plan_id=1072]

(8) Sort
Input [6]: [Payment Format#9, Transactions#96L, Laundering#97L, Laundering Rate (%)#114, Share of Laundering (%)#115, Relative Risk#116]
Arguments: [Laundering Rate (%)#114 DESC NULLS LAST], true, 0

(9) AdaptiveSparkPlan
Output [6]: [Payment Format#9, Transactions#96L, Laundering#97L, Laundering Rate (%)#114, Share of Laundering (%)#115, Relative Risk#116]
Arguments: isFinalPlan=false
\end{lstlisting}